% संपूर्ण मराठी ग्रंथातील प्रकरण 27: अपूर्णतेचा परिचय
% हा संपादनयोग्य स्रोतखंड आहे. संकलनासाठी संपूर्ण स्रोतसंग्रहातील
% अक्षररूपे, पूर्वभाग, चिन्हव्याख्या आणि आकृती-स्रोत आवश्यक आहेत.
% मूळ: Open Logic Project; मजकूर आणि रूपांतर CC BY 4.0.
% यंत्रानुवाद, दुरुस्ती आणि तपासणी: OpenAI Codex —
% GPT-5.6 Sol आणि GPT-6 Sol, दोन्ही Ultra effort.
% दुरुस्ती-पुनरावलोकन: GPT-6.1 Sol, Ultra effort.
\chapter{अपूर्णतेचा परिचय}\label{inc:int::chap}
\begin{editorial}
 या भागातील साहित्य अपूर्णतेच्या प्रमेयांविषयी आहे. यासाठी
 प्रथम-क्रम तर्कशास्त्राच्या भागातील साहित्य (विशेषतः सिद्धता-पद्धती)
 आणि संगणनक्षमता भागातील पुनरावर्ती फलनांविषयीचे साहित्य आवश्यक आहे.
 हे Jeremy Avigad यांच्या टिपणांवर आधारित असून Richard Zach यांनी त्यांत
 सुधारणा केल्या आहेत.
\end{editorial}









\section{ऐतिहासिक पार्श्वभूमी}\label{inc:int:bgr:sec}

या विभागात अपूर्णतेच्या प्रमेयांची पार्श्वभूमी समजण्यासाठी उपयुक्त अशा
ऐतिहासिक घडामोडींची थोडक्यात चर्चा करू. विशेषतः, गणितीय तर्कशास्त्राच्या
इतिहासाचा अगदी स्थूल आढावा घेऊ आणि मग गणिताच्या पायाभरणीच्या इतिहासाविषयी
थोडे बोलू.

\begin{digress}
  ``गणितीय तर्कशास्त्र'' या संज्ञेला दोन अर्थ आहेत. ``गणितीय'' हा शब्द
  अभ्यासविषय दर्शवतो असे मानल्यास, ``गणिताचे तर्कशास्त्र'' म्हणजे गणितीय
  युक्तिवादाची तत्त्वे. तो अभ्यासाची पद्धत दर्शवतो असे मानल्यास,
  ``तर्कशास्त्राचे गणित'' म्हणजे युक्तिवादाच्या तत्त्वांचा गणितीय अभ्यास.
  पुढील विवेचनात गणितीय तर्कशास्त्र हे दोन्ही अर्थांनी, अनेकदा एकाच वेळी,
  येते.
\end{digress}

तर्कशास्त्राचा अभ्यास मूलतः ॲरिस्टॉटलपासून सुरू झाला. तो सुमारे
384--322 \textsc{bce} या काळात जगला. त्याच्या \emph{कॅटेगरीज}, \emph{प्रायर
  ॲनॅलिटिक्स} आणि \emph{पोस्टीरियर ॲनॅलिटिक्स} या ग्रंथांत वैज्ञानिक युक्तिवादाच्या
तत्त्वांचा पद्धतशीर अभ्यास आहे; त्यात संवाक्याचाही सखोल व
पद्धतशीर अभ्यास समाविष्ट आहे.

मध्ययुगभर विद्वत् तत्त्वज्ञानावर ॲरिस्टॉटलच्या तर्कशास्त्राचे वर्चस्व होते.
अठराव्या शतकातसुद्धा कांटने त्याचे तर्कशास्त्र परिपूर्ण असून त्यात सुधारणा
करण्याची गरज नाही, असे म्हटले. पण न्यायवाक्यरचनेचा सिद्धांत गणितीय
युक्तिवादाच्या अगदी वरवरच्या पैलूंपलीकडे काही दाखवण्यासाठी अत्यंत मर्यादित
आहे. त्याच्या एक शतक आधी, न्यूटनचा समकालीन लाइब्निझ याने तार्किक
युक्तिवादासाठी पूर्ण ``कलन'' असावे अशी कल्पना केली होती. त्याने अशा कलनाची
रचना करण्याच्या दिशेने काही प्राथमिक पावलेही टाकली; त्यातून मूलतः विधानीय
तर्कशास्त्राची एक आवृत्ती वर्णिली गेली.

एकोणिसावे शतक तर्कशास्त्रासाठी निर्णायक ठरले. 1854 मध्ये जॉर्ज बूलने
\emph{विचाराचे नियम} लिहिला; त्यातील विधानीय तर्कशास्त्राचा सखोल बीजगणितीय
अभ्यास आधुनिक मांडण्यांपासून फारसा दूर नाही. 1879 मध्ये गॉटलोब फ्रेगेने
\emph{Begriffsschrift} (संकल्पनालेखन) प्रसिद्ध केला. त्यात विधानीय
तर्कशास्त्राला संख्यापक आणि संबंध यांची भर घातली असल्याने प्रथम-क्रम
तर्कशास्त्र समाविष्ट होते. प्रत्यक्षात, फ्रेगेच्या तार्किक पद्धतींत
उच्च-क्रम तर्कशास्त्र आणि त्याहून अधिक गोष्टीही होत्या. \emph{अंकगणिताचे
  मूलभूत नियम} या ग्रंथात त्याने सर्व अंकगणित केवळ तार्किक गृहीतकांपासून
आपल्या Begriffsschrift मध्ये निष्पन्न करता येईल, हे दाखवण्याचा प्रयत्न
केला. दुर्दैवाने, रसेलने 1902 मध्ये दाखवून दिल्याप्रमाणे ही गृहीतके
विसंगत निघाली. तरी तो विसंगत स्वयंसिद्धक बाजूला ठेवल्यास, फ्रेगेने आधुनिक
तर्कशास्त्र जवळजवळ एकट्याने निर्माण केले; ही विलक्षण कामगिरी होती.
बूलनंतर पिअर्स आणि श्रॉडर यांच्यासह बीजगणितीय दृष्टिकोन असलेल्या इतर
विचारवंतांनीही संख्यापकाधारित तर्कशास्त्र स्वतंत्रपणे विकसित केले.

आता गणिताच्या पायाभरणीतील घडामोडींकडे वळू. तर्कशास्त्राची गणितात महत्त्वाची
भूमिका असल्यामुळे नुकत्याच वर्णिलेल्या घडामोडींशी त्यांचा मोठा संबंध आहे.
उदाहरणार्थ, सर्व अंकगणित केवळ आपल्या तार्किक चौकटीवर आधारता येईल हे
दाखवण्याच्या स्पष्ट उद्देशाने फ्रेगेने आपले तर्कशास्त्र विकसित केले.
विशेषतः, कांटने अंकगणितातील सत्ये अनुभवपूर्व \emph{संश्लेषक} मानली होती;
फ्रेगेला त्याऐवजी ती अनुभवपूर्व \emph{विश्लेषक} आहेत, हे दाखवायचे होते.
भूमितीतील सत्ये अनुभवपूर्व संश्लेषक आहेत हे मात्र फ्रेगेने मान्य केले.


अनेकांच्या मते गणिताचा खरा उदय ग्रीकांबरोबर झाला. सुमारे 300 इ.स.पू.
मध्ये लिहिलेला युक्लिडचा \emph{मूलतत्त्वे} हा ग्रंथ काटेकोरपणा आणि
अचूकता यांवर भर देणाऱ्या ग्रीक गणिताचे विकसित उदाहरण आहे. युक्लिडच्या
\emph{मूलतत्त्वे}मधील व्याख्या आणि सिद्धता आजही माध्यमिक शाळेच्या
भूमितीच्या पाठ्यपुस्तकांत जवळजवळ तशाच टिकून आहेत (जिथे अजून माध्यमिक
शाळेत भूमिती शिकवली जाते तिथे). गणितातच नव्हे तर तत्त्वज्ञानाच्या
शाखांतही काटेकोर युक्तिवादाचा आदर्श म्हणून गणितीय युक्तिवादाचे हे रूप
मानले गेले आहे. (स्पिनोझाने नैतिक आणि धार्मिक युक्तिवादसुद्धा
युक्लिडच्या शैलीत मांडले; ते पाहणे काहीसे विचित्र वाटते!)

न्यूटन आणि लाइब्निझ यांनी सतराव्या शतकात कलनाचा शोध लावला. (प्राधान्यावरून
शतकानुशतके तीव्र वाद चालला; पण आज बहुतेक अभ्यासकांच्या मते त्यांची
बहुतांश प्रगती स्वतंत्रपणे झाली.) कलनात, उदाहरणार्थ, अनंतसूक्ष्म
राशींच्या अनंत बेरजांबद्दल युक्तिवाद करावा लागतो. या वैशिष्ट्यांमुळे
बिशप बर्कलीने टीका केली: ईश्वरावरील विश्वास आपल्या काळातील गणितापेक्षा
कमी तर्कसंगत नाही, असे त्याचे म्हणणे होते. अठराव्या शतकात कलनाच्या
पद्धती मोठ्या प्रमाणावर वापरल्या गेल्या; उदाहरणार्थ, लिओनहार्ड
ऑयलरने अनंत बेरजांचा समावेश असलेली गणिते करून प्रभावी निष्कर्ष मिळवले.

एकोणिसाव्या शतकात गणितज्ञांनी कलनाला अधिक भक्कम पाया देऊन बर्कलीच्या
टीकेला उत्तर देण्याचा प्रयत्न केला. कॉशी, वायेरश्ट्रास, बोल्झानो आणि
इतरांच्या प्रयत्नांतून ``एप्सिलॉन आणि डेल्टा'' यांच्या साहाय्याने सीमा,
सातत्य, अवकलन आणि समाकलन यांच्या आजच्या व्याख्या निर्माण झाल्या; म्हणजेच
त्यात अनंतसूक्ष्म राशींचा संदर्भ उरला नाही. पुढे त्याच शतकात गणितज्ञांनी
वास्तव संख्यांसह कलनाचे सर्व पैलू नैसर्गिक संख्यांच्या आधारे स्पष्ट
करण्याचा अधिक प्रयत्न केला. (क्रोनेकर: ``पूर्ण संख्या ईश्वराने निर्माण
केल्या; उरलेले सर्व मानवाचे कार्य आहे.'') 1872 मध्ये डेडेकिंडने
``सातत्य आणि अपरिमेय संख्या'' लिहून वास्तव संख्या परिमेय संख्यांच्या
संचांच्या रूपात कशा ``रचता'' येतात हे दाखवले. (परिमेय संख्यांकडे, तुम्हाला
माहीत आहेच, नैसर्गिक संख्यांच्या जोड्यांच्या रूपात पाहता येते.) 1888 मध्ये
त्याने ``Was sind und was sollen die Zahlen'' (साधारणतः ``नैसर्गिक
संख्या म्हणजे काय आणि त्यांचे स्वरूप कसे असावे?'') लिहून नैसर्गिक संख्या
निखळ ``तार्किक'' संज्ञांत स्पष्ट करण्याचा प्रयत्न केला. 1887 मध्ये
क्रोनेकरने ``\"Uber den Zahlbegriff'' (``संख्येच्या संकल्पनेविषयी'')
लिहून सर्व गणितीय वस्तू पूर्ण संख्यांच्या आधारे दर्शवण्याविषयी सांगितले.
1889 मध्ये ज्यूझेप्पे पेआनोने नैसर्गिक संख्यांसाठी आकारिक, प्रतीकात्मक
स्वयंसिद्धके दिली.

एकोणिसाव्या शतकाच्या अखेरीस अनंताशी व्यवहार करण्याचे नवे धाडसही दिसले.
त्याआधी अनंत वस्तू आणि रचना (उदा., नैसर्गिक संख्यांचा संच) अतिशय
सावधपणे हाताळल्या जात. ``अनंत अनेक'' म्हणजे ``तुम्हाला हव्या तितक्या''
आणि ``सीमेत जवळ जाणे'' म्हणजे ``तुम्हाला हवे तितके जवळ जाणे'' असे समजले
जाई. पण गेऑर्ग कांटोरने अनंताचा प्रत्यक्ष अर्थाने विचार करणे शक्य आहे हे
दाखवले. कांटोर, डेडेकिंड आणि इतरांच्या कार्यामुळे गणिताची आज व्यापकपणे
स्वीकारली जाणारी सर्वसाधारण संचसिद्धान्तीय समज विकसित होऊ लागली.

यातून आपण तर्कशास्त्र आणि पायाभरणीतील विसाव्या शतकाच्या घडामोडींपर्यंत
पोहोचतो. 1902 मध्ये रसेलने फ्रेगेला पत्र लिहून त्याच्या तार्किक पद्धतीतील
विरोधाभास कळवला.

1904 मध्ये झर्मेलोने तथाकथित ``निवडीचे स्वयंसिद्धक'' वापरून कांटोरचे
सुव्यवस्थापन तत्त्व सिद्ध केले; या स्वयंसिद्धकाच्या स्वीकारार्हतेवर मोठा वाद
झाला. 1910 ते 1913 या काळात रसेल आणि व्हाइटहेड यांच्या \emph{Principia
  Mathematica}चे तीन खंड प्रसिद्ध झाले. गणिताला तार्किक आधार देण्याचा
फ्रेगेचा कार्यक्रम त्यांनी पुढे नेला. मात्र रसेल आणि व्हाइटहेड यांना
निखळ तार्किक म्हणून समर्थनीय वाटणे कठीण अशी दोन तत्त्वे स्वीकारावी
लागली: अनंततेचे स्वयंसिद्धक आणि ``न्यूनीकरणीयतेचे'' स्वयंसिद्धक. 1900
च्या दशकात पॉँकारेने गणितातील ``स्वसमावेशक व्याख्यां''च्या वापरावर
टीका केली. 1910 च्या दशकात ब्रौवरने विमध्य सिद्धांत
($\varphi \lor \lnot \varphi$) न वापरता संपूर्ण गणिताची ``अंतःप्रज्ञावादी''
पायावर पुनर्मांडणी करण्याचा प्रस्ताव मांडायला सुरुवात केली.

खरोखरच विलक्षण काळ!{} सर्व गणित तर्कशास्त्रात रूपांतरित करण्याच्या
कार्यक्रमाला आता ``तर्कमीमांसावाद'' म्हणतात. वर सांगितलेल्या अडचणींमुळे
तो अयशस्वी ठरला, असे सामान्यतः मानले जाते. मानसिक अंतःप्रज्ञावादी
रचनांवर गणित उभारण्याच्या कार्यक्रमाला ``अंतःप्रज्ञावाद'' म्हणतात;
दैनंदिन गणितावर तो अति कठोर निर्बंध घालतो, असे मानले जाते. शतकाच्या
वळणावर, सर्वकाळातील अत्यंत प्रभावशाली गणितज्ञांपैकी एक असलेला डेव्हिड
हिल्बर्ट हा कांटोर आणि डेडेकिंड यांनी आणलेल्या नव्या, अमूर्त पद्धतींचा
ठाम समर्थक होता: ``कांटोरने आपल्यासाठी निर्माण केलेल्या स्वर्गातून
आपल्याला कोणीही हाकलू शकणार नाही.'' त्याच वेळी, या नव्या पद्धतींवरील
पायाभूत टीकेचीही त्याला जाणीव होती (विचित्र म्हणजे, आता या पद्धतींना
``अभिजात'' म्हणतात). दोन्ही गोष्टी साधण्याचा त्याने असा मार्ग सुचवला:
\begin{enumerate}
\item अभिजात पद्धती आकारिक स्वयंसिद्धके आणि नियम यांच्या साहाय्याने
  दर्शवा; गणितीय प्रश्न स्वयंसिद्धकीय पद्धतीतील सूत्रे म्हणून दर्शवा.
\item या आकारिक निगमन-पद्धती सुसंगत आहेत हे सिद्ध करण्यासाठी सुरक्षित,
  ``सांतवादी'' पद्धती वापरा.
\end{enumerate}

पहिले उद्दिष्ट साधण्याच्या दिशेने हिल्बर्टच्या कार्याने मोठी मजल मारली.
1899 मध्ये त्याने आपल्या प्रसिद्ध \emph{भूमितीची पायाभरणी} या पुस्तकात
भूमितीसाठी हे काम केले होते. पुढील वर्षांत, हिल्बर्टने भूमितीसाठी जे
केले ते गणिताच्या इतर शाखांत करण्यासाठी तो, त्याचे अनेक विद्यार्थी आणि
सहकारी यांनी काम केले. हिल्बर्टने स्वतः अंकगणित आणि विश्लेषण यांसाठी
स्वयंसिद्धक-प्रणाली दिल्या. झर्मेलोने संचसिद्धान्ताची स्वयंसिद्धकीय मांडणी
दिली; फ्रेंकेल, स्कोलेम, फोन न्यूमन आणि इतरांनी ती पुढे विकसित केली.
1920 च्या दशकाच्या मध्यापर्यंत ``सर्व'' गणिताची स्वयंसिद्धकीय मांडणी
म्हणवणारे दोन दृष्टिकोन होते: रसेल आणि व्हाइटहेड यांचे \emph{Principia
  mathematica}, आणि पुढे झर्मेलो--फ्रेंकेल संच उपपत्ती म्हणून ओळखली
जाणारी मांडणी.

1921 मध्ये हिल्बर्टने या पद्धती सुसंगत असल्याचे सिद्ध करण्याच्या उद्देशाने
संशोधन कार्यक्रम सुरू केला. यात त्याला अनेक विद्यार्थ्यांची, विशेषतः
बेर्नाइस, आकरमान आणि पुढे गेन्ट्सेन यांची मदत झाली. या उद्दिष्टामागील
मूळ कल्पना अशी होती: गणितात विसंगतीची निष्पत्ती शक्य आहे का, हा
प्रश्न चिन्हांच्या संभाव्य क्रमांबद्दलच्या संयुक्तिक प्रश्नात रूपांतरित
करायचा. म्हणजे, अंकगणित, विश्लेषण किंवा संचसिद्धान्त यांच्या
स्वयंसिद्धक-प्रणालीतील स्वयंसिद्धकांपासून, उदाहरणार्थ, $\varphi \land
\lnot \varphi$ ची योग्य निष्पत्ती ठरणाऱ्या वाक्यांचे संभाव्य क्रम
तपासायचे. असा चिन्हक्रम अशक्य असल्याची सिद्धता ही स्वतः गणितीय सिद्धता
असल्यामुळे या स्वयंसिद्धकीय पद्धतींत आकारिकरीत्या मांडता आली असती.
दुसऱ्या शब्दांत, उदाहरणार्थ, अंकगणित सुसंगत आहे असे सांगणारे
$\OCon$ हे वाक्य असले असते. शिवाय, विशेषतः त्याची सिद्धता केवळ अतिशय
मर्यादित, ``सांतवादी'' साधनांनी होत असल्यास, ते वाक्य संबंधित पद्धतींत
सिद्ध करता आले पाहिजे होते.

विकसित स्वयंसिद्धक-प्रणाली प्रत्येक गणितीय प्रश्न सोडवतील, हे दुसरे
उद्दिष्ट दोन प्रकारे नेमके करता येते. पहिला प्रकार असा: गणिताच्या
स्वयंसिद्धक-प्रणालीच्या भाषेतील कोणतेही वाक्य~$\varphi$ घेतल्यास,
स्वयंसिद्धकांपासून $\varphi$ किंवा $\lnot \varphi$ सिद्ध करता येते. असे असल्यास,
स्वयंसिद्धकांच्या आधारे ज्यांची सिद्धता किंवा खंडन दोन्ही होत नाहीत अशी
वाक्ये, म्हणजे त्यांनी न सोडवलेले प्रश्न, उरणार नाहीत. हा गुणधर्म
असलेल्या स्वयंसिद्धक-प्रणालीला \emph{संपूर्ण} म्हणतात. अर्थात, दिलेल्या
एखाद्या वाक्यासाठी या दोन पर्यायांपैकी कोणता खरा आहे हे ठरवणे तरीही
कठीण असू शकते. पण तत्त्वतः ते ठरवण्याची पद्धत असली पाहिजे. प्रत्यक्षात,
हिल्बर्टने विचारात घेतलेल्या स्वयंसिद्धक आणि निष्पत्ती-पद्धतींसाठी
संपूर्णतेवरून अशी पद्धत अस्तित्वात आहे असे निष्पन्न होईल---जरी हिल्बर्टला
याची जाणीव नव्हती. दुसरा अर्थ अधिक प्रबळ मागणी करतो: दिलेले
वाक्य~$\varphi$ स्वयंसिद्धकांपासून निष्पन्न करता येते की नाही हे ठरवणारी
यांत्रिक, संगणकीय पद्धत असावी.

1931 मध्ये गोडेलने दोन ``अपूर्णतेची प्रमेये'' सिद्ध केली. त्यांवरून
हा कार्यक्रम मूळ स्वरूपात यशस्वी होऊ शकत नाही, हे दिसले. अंकगणितासाठी
पुरेशी प्रबळ, सुसंगत आणि प्रभावी स्वयंसिद्धकीय मांडणी असलेली प्रणाली
संपूर्ण असू शकत नाही. शिवाय, दुसऱ्या अपूर्णता प्रमेयाच्या अटी पूर्ण करणारी
अशी प्रणाली तिच्या प्रमाण सिद्धतायोग्यता-विधेयाने व्यक्त केलेली स्वतःची
सुसंगतता सिद्ध करू शकत नाही. त्या सुसंगतता-विधानाचे खंडनही होत नाही असे
म्हणण्यासाठी अधिक गृहीतक हवे; उदाहरणार्थ, प्रणालीने सिद्ध केलेली सर्व
अंकगणितीय वाक्ये नैसर्गिक संख्यांच्या प्रमाण प्रतिमानात सत्य असावीत.


यामुळे हिल्बर्टच्या मूळ कार्यक्रमाला घातक धक्का बसला. तरीही गणितात
अनेकदा घडते तसे, यातून संशोधनाच्या नव्या दिशाही उघडल्या. सर्व गणिताला
व्यापणारी एकच आकारिक पद्धत नसल्यास, अधिक मर्यादित पद्धती विकसित करून
त्यांत काय सिद्ध करता येते याचा अभ्यास करणे अर्थपूर्ण ठरते. या पद्धतींची
सुसंगतता सिद्ध करण्यासाठी कमी निर्बंध असलेल्या सिद्धता-पद्धती विकसित
करणे, आणि त्यांची सुसंगतता सिद्ध करणे किती कठीण आहे हे मोजण्याचे मार्ग
शोधणेही अर्थपूर्ण ठरते. अपूर्णतेच्या प्रमेयांची गृहीतके पूर्ण करणाऱ्या
आकारिक पद्धतींत असे प्रश्न असतात की जे त्या सोडवू शकत नाहीत. म्हणून
दिलेल्या आकारिक पद्धतीला न सोडवता येणारे ``महत्त्वाचे'' प्रश्न शोधणे,
आणि ते सोडवण्यासाठी पद्धत किती प्रबळ असावी लागते हे ठरवणे योग्य ठरते.
आजतागायत तर्कशास्त्रज्ञ या प्रश्नांचा सिद्धता-उपपत्ती या नव्या
गणितीय शाखेत अभ्यास करत आहेत.












\section{व्याख्या}\label{inc:int:def:sec}

गणिताला आकारिक रूप देण्याचा आणि असे रूप सुसंगत व संपूर्ण आहे हे
दाखवण्याचा हिल्बर्टचा कार्यक्रम राबवण्यासाठी प्रथम भाषा, तार्किक चौकट
आणि स्वयंसिद्धक-प्रणाली निवडावी लागेल. आपल्या उद्देशासाठी गणित प्रथम-क्रम
भाषेत आकारिकरीत्या मांडता येते असे समजू. म्हणजे स्थिरांक,
फलने आणि विधेये यांचा असा काही संच आहे की प्रथम-क्रम
तर्कशास्त्रातील संयोजके व संख्यापक यांच्यासह त्यातून गणितातील विधाने
व्यक्त करता येतात. अशी भाषा अस्तित्वात आहे याबद्दल बहुतेकांचे एकमत
आहे: ती संचसिद्धान्ताची भाषा आहे आणि त्यात $\in$ हे एकमेव अतार्किक
चिन्ह आहे. इतकी साधी भाषा इतकी अभिव्यक्तिक्षम असू शकते हा दावा
पहिल्या नजरेला नक्कीच अविश्वसनीय वाटतो; गणितातील जवळजवळ सर्व काही
या अत्यल्प शब्दसंग्रहात व्यक्त करता येते हे प्रस्थापित करण्यासाठी
मोठे काम करावे लागले. सध्या चर्चा सोपी ठेवण्यासाठी ती केवळ
अंकगणितापुरती, म्हणजे नैसर्गिक संख्या~$\Nat$ यांच्याशी संबंधित
गणिताच्या भागापुरती, मर्यादित करू. अंकगणितातील तथ्ये व्यक्त करण्यासाठी
योग्य भाषा~$\Lang L_A$ आहे. $\Lang L_A$ मध्ये एक द्विस्थानी
विधेय~$<$, एक स्थिरांक~$\Obj 0$, एक एकस्थानी
फलन~$\prime$, आणि $+$ व $\times$ ही दोन द्विस्थानी
फलने आहेत.

\begin{defn}
वाक्यांचा संच~$\Gamma$ हा तार्किक निष्पन्नतेखाली बंद असेल, म्हणजे
$\Gamma = \Setabs{\varphi}{\Gamma \Entails
\varphi}$ असेल, तर त्याला \emph{उपपत्ती} म्हणतात.
\end{defn}

उपपत्ती निर्दिष्ट करण्याचे दोन सोपे मार्ग आहेत. एक म्हणजे एखाद्या
संरचनांमध्ये सत्य असलेल्या वाक्यांचा संच म्हणून
तिला निर्दिष्ट करणे. उदाहरणार्थ, $\Lang L_A$ साठी अशी संरचना
घ्या की तिचे प्रांत~$\Nat$ आहे आणि सर्व अतार्किक चिन्हांचा
अपेक्षेप्रमाणे अर्थ लावला आहे.

\begin{defn}\label{inc:int:def:def:standard-model}
\emph{अंकगणिताचे प्रमाण प्रतिमान} ही पुढीलप्रमाणे परिभाषित
केलेली संरचना~$\Struct
N$ आहे:
\begin{enumerate}
\item $\Domain N = \Nat$
\item $\Assign{\Obj 0}{N} = 0$
\item $\Assign{\Obj \prime}{N}(n) = n + 1$ सर्व $n \in \Nat$ साठी
\item $\Assign{\Obj +}{N}(n, m) = n + m$ सर्व $n, m \in \Nat$ साठी
\item $\Assign{\Obj \times}{N}(n, m) = n\cdot m$ सर्व $n, m \in \Nat$ साठी
\item $\Assign{\Obj <}{N} = \Setabs{\tuple{n, m}}{n \in \Nat, m \in
  \Nat, n < m}$
\end{enumerate}
\end{defn}

`$\times$' आणि `$\cdot$' यांतील फरक लक्षात घ्या: $\times$ हे
अंकगणिताच्या भाषेतील चिन्ह आहे. गुणाकाराची आठवण व्हावी म्हणून ते
निवडले आहे; पण $\times$ ही गुणाकाराची क्रिया नसून द्विस्थानी फलनचिन्ह
आहे (औपचारिकरीत्या, $\Obj f^2_1$). याउलट, $\cdot$ हेच नेहमीचे
गुणाकाराचे फलन \emph{आहे}. $n \cdot m$ सारखे काही दिसल्यास त्याचा
अर्थ $n$ आणि~$m$ या संख्यांचा गुणाकार होतो; $x \times y$ सारखे
काही दिसल्यास आपण अंकगणिताच्या भाषेतील पदाबद्दल बोलत असतो. प्रमाण
प्रतिमानात फलनचिन्ह~$\times$ चा अर्थ नैसर्गिक संख्यांवरील फलन
$\cdot$ असा निश्चित केला जातो. बेरीजेसाठी मात्र अंकगणिताच्या
भाषेतील फलनचिन्ह आणि नैसर्गिक संख्यांवरील बेरीज हे दोन्ही दर्शवण्यासाठी
आपण $+$ वापरतो. त्याचा अर्थ संदर्भावरून ठरवावा लागतो.

\begin{defn}
\emph{सत्य अंकगणिताची} उपपत्ती म्हणजे अंकगणिताच्या प्रमाण प्रतिमानात
पूर्त होणाऱ्या वाक्यांचा संच; म्हणजे,
\[\Th{TA} = \Setabs{\varphi}{\Sat{N}{\varphi}}.
\]
\end{defn}

$\Th{TA}$ ही उपपत्ती आहे. कारण $\Th{TA} \Entails \varphi$ असेल तेव्हा
$\Th{TA}$ ची पूर्ती करणाऱ्या प्रत्येक संरचनांमध्ये $\varphi$ ची
पूर्ती होते. $\Sat{N}{\Th{TA}}$ असल्यामुळे $\Sat{N}{\varphi}$ मिळते;
म्हणून $\varphi \in \Th{TA}$.

उपपत्ती~$\Gamma$ निर्दिष्ट करण्याचा दुसरा मार्ग म्हणजे काही
वाक्यसंच~$\Gamma_0$ पासून तार्किकरीत्या निष्पन्न होणाऱ्या
वाक्यांचा संच म्हणून तिला सांगणे. तेव्हा $\Gamma$ हे
तार्किक निष्पन्नतेखाली $\Gamma_0$ चे ``संवरण'' असते. या मार्गाने उपपत्ती निर्दिष्ट करणे तेव्हाच
उपयुक्त ठरते जेव्हा $\Gamma_0$ स्पष्टपणे निर्दिष्ट केला असेल;
उदाहरणार्थ, $\Gamma_0$ मधील घटक ची यादी दिली असेल.
किमान एवढे तरी हवे की $\Gamma_0$ निर्णेय असावा: एखादे
वाक्य $\Gamma_0$ चे घटक आहे की नाही हे ठरवण्यासाठी
संगणनक्षम चाचणी असावी. $\Gamma_0$ मधील वाक्यांना
$\Gamma$ ची \emph{स्वयंसिद्धके}, आणि $\Gamma$ ला
$\Gamma_0$ या संचातील \emph{स्वयंसिद्धकांद्वारे निरूपित} म्हणतो.

\begin{defn}
उपपत्ती~$\Gamma$ ही $\Gamma_0$ या संचातील
\emph{स्वयंसिद्धकांद्वारे निरूपित} आहे तेव्हा आणि केवळ तेव्हाच
\[\Gamma = \Setabs{\varphi}{\Gamma_0 \Entails \varphi}
\]
असते.
\end{defn}

\begin{defn}
पुढील वाक्यांनी स्वयंसिद्धकांद्वारे निरूपित केलेली उपपत्ती
$\Th{Q}$ ``रॉबिन्सनची $\Th{Q}$'' म्हणून ओळखली जाते आणि ती
अंकगणिताची अतिशय साधी उपपत्ती आहे.
\begin{align*}
& \lforall[x][\lforall[y][(\eq[x'][y'] \lif \eq[x][y])]] \tag{$\mathsf{Q}_1$}\\
& \lforall[x][\eq/[\Obj 0][x']] \tag{$\mathsf{Q}_2$}\\
& \lforall[x][(\eq[x][\Obj 0] \lor \lexists[y][\eq[x][y']])] \tag{$\mathsf{Q}_3$}\\
& \lforall[x][\eq[(x + \Obj 0)][x]] \tag{$\mathsf{Q}_4$}\\
& \lforall[x][\lforall[y][\eq[(x + y')][(x + y)']]] \tag{$\mathsf{Q}_5$}\\
& \lforall[x][\eq[(x \times \Obj 0)][\Obj 0]] \tag{$\mathsf{Q}_6$}\\
& \lforall[x][\lforall[y][\eq[(x \times y')][((x \times y) + x)]]] \tag{$\mathsf{Q}_7$}\\
& \lforall[x][\lforall[y][(x < y \liff \lexists[z][\eq[(z' + x)][y]])]] \tag{$\mathsf{Q}_8$}
\end{align*}
$\{\mathsf{Q}_1, \dots, \mathsf{Q}_8\}$ हा वाक्यांचा संच $\Th{Q}$ ची
स्वयंसिद्धके आहेत. म्हणून त्यांच्यापासून तार्किकरीत्या निष्पन्न होणारी सर्व
वाक्ये $\Th{Q}$ मध्ये आहेत:
\[\Th{Q} = \Setabs{\varphi}{\{\mathsf{Q}_1, \dots, \mathsf{Q}_8\} \Entails \varphi}.
\]
\end{defn}

\begin{defn}
$\varphi(x)$ हे $\Lang L_A$ मधील असे सूत्र आहे असे समजा की
त्यात $x$ आणि $y_1$, \dots, $y_n$ ही मुक्त चरे आहेत. तर पुढील
रूपातील कोणतेही वाक्य हे \emph{विगमन-रूपबंधाचे} उदाहरण आहे:
\[\lforall[y_1][\dots\lforall[y_n][((\varphi(\Obj 0) \land \lforall[x][(\varphi(x)
\lif \varphi(x'))]) \lif \lforall[x][\varphi(x)])]]
\]

\emph{पेआनो अंकगणित}~$\Th{PA}$ ही $\Th{Q}$ च्या स्वयंसिद्धकांसह
विगमन-रूपबंधाच्या सर्व उदाहरणांनी स्वयंसिद्धकांद्वारे निरूपित
केलेली उपपत्ती आहे.
\end{defn}

\begin{explain}
विगमन-रूपबंधाचे प्रत्येक उदाहरण~$\Struct{N}$ मध्ये सत्य असते.
सूत्र~$\varphi$ मध्ये फक्त एकच मुक्त चल~$x$ असल्यास
हे सर्वात सहज दिसते. तेव्हा $\varphi(x)$ हे~$\Struct{N}$ मध्ये
$\Nat$ चा उपसंच~$X_{\varphi}$ परिभाषित करते. $s(x) = n$ असताना
$\Sat{N}{\varphi(x)}[s]$ होईल अशा सर्व~$n \in \Nat$ चा~$X_{\varphi}$ हा संच आहे.
विगमन-रूपबंधाचे संबंधित उदाहरण असे आहे:
\[((\varphi(\Obj 0) \land \lforall[x][(\varphi(x) \lif \varphi(x'))]) \lif
  \lforall[x][\varphi(x)]).
\]
याचा पूर्वपक्ष~$\Struct{N}$ मध्ये सत्य असल्यास $0 \in X_{\varphi}$;
आणि $n \in X_{\varphi}$ असेल तेव्हा $n+1$ देखील त्यात असते.
$0 \in X_{\varphi}$ असल्याने $1 \in X_{\varphi}$ मिळते. पुन्हा
$1 \in X_{\varphi}$ असल्याने $2 \in X_{\varphi}$ मिळते, आणि पुढे असेच.
म्हणून प्रत्येक
$n \in \Nat$ साठी $n \in X_{\varphi}$. याचा अर्थ
$\lforall[x][\varphi(x)]$ ची~$\Struct{N}$ मध्ये पूर्ती होते.
\end{explain}

$\Th{Q}$ आणि $\Th{PA}$ या दोन्ही स्वयंसिद्धकांद्वारे निरूपित
उपपत्ती आहेत. त्या किती प्रबळ आहेत, हा मुख्य प्रश्न आहे.
उदाहरणार्थ, $\Lang L_A$ मध्ये व्यक्त करता येणारी~$\Nat$ विषयीची
सर्व सत्ये $\Th{PA}$ सिद्ध करू शकते का? नेमके सांगायचे तर,
$\Lang L_A$ मध्ये मांडता येणारे सर्व प्रश्न $\Th{PA}$ ची
स्वयंसिद्धके सोडवतात का?

हा प्रश्न दुसऱ्या रीतीनेही विचारता येईल: $\Th{PA} = \Th{TA}$
आहे का? $\Th{TA}$ अर्थातच~$\Nat$ विषयीची सर्व सत्ये सिद्ध
करते (म्हणजे ती त्यात समाविष्ट असतात), आणि $\Lang L_A$ मधील
सर्व प्रश्न सोडवते. कारण $\varphi$ हे $\Lang L_A$ मधील वाक्य
असेल तर $\Sat{N}{\varphi}$ किंवा $\Sat{N}{\lnot \varphi}$ यांपैकी
एक सत्य असते; म्हणून $\Th{TA}
\Entails \varphi$ किंवा $\Th{TA} \Entails \lnot \varphi$. अशा उपपत्तीला
\emph{संपूर्ण} म्हणतात.

\begin{defn}
उपपत्ती $\Gamma$ ही \emph{संपूर्ण} आहे तेव्हा आणि केवळ तेव्हाच तिच्या
भाषेतील प्रत्येक वाक्य~$\varphi$ साठी $\Gamma \Entails \varphi$
किंवा $\Gamma \Entails \lnot \varphi$.
\end{defn}

\begin{explain}
संपूर्णता प्रमेयामुळे $\Gamma \Entails \varphi$ तेव्हा आणि केवळ तेव्हाच
$\Gamma \Proves
\varphi$. म्हणून $\Gamma$ संपूर्ण आहे तेव्हा आणि केवळ तेव्हाच तिच्या भाषेतील
प्रत्येक वाक्य~$\varphi$ साठी $\Gamma \Proves \varphi$ किंवा
$\Gamma \Proves \lnot \varphi$.
\end{explain}

आणखी एक प्रश्न असा आहे: एखादे वाक्य~$\Th{TA}$ मध्ये,
$\Th{PA}$ मध्ये, किंवा अगदी~$\Th{Q}$ मध्येही आहे की नाही हे
तपासण्यासाठी संगणकीय प्रक्रिया वापरता येईल का? संच (उदाहरणार्थ,
वाक्यांचा संच) निर्णेय कधी असतो हे परिभाषित करून हा
प्रश्न अधिक नेमका करता येईल.

\begin{defn}
संच $X$ हा \emph{निर्णेय} आहे तेव्हा आणि केवळ तेव्हाच अशी संगणकीय प्रक्रिया
असते की आदान~$x$ दिल्यावर $x \in X$ असल्यास ती~$1$ आणि
अन्यथा~$0$ परत करते.
\end{defn}

म्हणून प्रश्न असा होतो: $\Th{TA}$ ($\Th{PA}$, $\Th{Q}$) निर्णेय आहे का?

या सर्व प्रश्नांचे उत्तर नकारार्थी असेल. यांपैकी कोणतीही उपपत्ती
निर्णेय नाही. मात्र ही घटना केवळ या उपपत्तींनाच लागू नाही. प्रत्यक्षात,
विशिष्ट अटी पूर्ण करणाऱ्या \emph{कोणत्याही} उपपत्तीसाठी हेच निष्कर्ष
लागू होतात. यांपैकी एक अट, जी $\Th{Q}$ आणि $\Th{PA}$ पूर्ण करतात,
अशी आहे की उपपत्ती निर्णेय स्वयंसिद्धकसंचाने निरूपित झालेली असावी.

\begin{defn}
एखादी उपपत्ती निर्णेय स्वयंसिद्धकसंचाने निरूपित झाली असेल तर ती
\emph{निर्णेय स्वयंसिद्धकसंचाने निरूपणयोग्य} आहे.
\end{defn}

\begin{ex}
सांत वाक्यांच्या संचाने स्वयंसिद्धकांद्वारे निरूपित
झालेली कोणतीही उपपत्ती निर्णेय स्वयंसिद्धकसंचाने निरूपणयोग्य असते, कारण प्रत्येक
सांत संच निर्णेय असतो. म्हणून, उदाहरणार्थ, $\Th{Q}$ ही
निर्णेय स्वयंसिद्धकसंचाने निरूपणयोग्य आहे.

$\Th{PA}$ सारखी रूपबंधांद्वारे स्वयंसिद्धके दिलेली उपपत्तीदेखील
निर्णेय स्वयंसिद्धकसंचाने निरूपणयोग्य असते. कारण $\psi$ हे $\Th{PA}$ च्या स्वयंसिद्धकांपैकी
आहे का, म्हणजे $\psi$ हे $\Th{PA}$ चे स्वयंसिद्धक असल्यास
$\Char{X}(\psi) = 1$ आणि
अन्यथा $= 0$ असणारे फलन $\Char{X}$ कसे काढायचे, हे पुढीलप्रमाणे
तपासता येते. प्रथम, $\psi$ हे $\Th{Q}$ च्या स्वयंसिद्धकांपैकी आहे
का ते तपासा. असल्यास उत्तर ``होय'' आणि $\Char{X}(\psi) = 1$.
नसल्यास ते विगमन-रूपबंधाचे उदाहरण आहे का हे तपासा. हे पद्धतशीरपणे
करता येते. येथे पुढील पद्धत कदाचित सर्वात स्पष्ट आहे: विगमन-रूपबंधाचे
प्रत्येक उदाहरण सुरुवातीला काही वैश्विक संख्यापकांनी, त्यानंतर
सोपाधिक असलेल्या एका उप-सूत्र ने सुरू होते. त्या सोपाधिकाचा
उत्तरपक्ष $\lforall[x][\varphi(x, y_1, \dots, y_n)]$ असतो. येथे $x$
आणि $y_1$, \dots, $y_n$ ही~$\varphi$ ची सर्व मुक्त चरे असतात, आणि
$\psi$ च्या सुरुवातीचे संख्यापक $y_1$, \dots,~$y_n$ ही चरे बद्ध
करतात. हे~$\varphi$ काढून घेतल्यावर त्याची मुक्त चरे सुरुवातीच्या
वैश्विक संख्यापकांनी आणि~$\lforall[x]$ ने बद्ध केलेल्या चरांशी
जुळतात का ते तपासा. मग सोपाधिकाचा पूर्वपक्ष पुढील स्वरूपाशी जुळतो
का ते तपासा:
\[\varphi(\Obj 0, y_1, \dots, y_n) \land \lforall[x][(\varphi(x, y_1, \dots, y_n)
\lif \varphi(x', y_1, \dots, y_n))]
\]
तो जुळल्यास $\psi$ विगमन-रूपबंधाचे उदाहरण आहे; न जुळल्यास $\psi$ तसे नाही.
\end{ex}

या प्रश्नाचे---आणि कोणत्या उपपत्ती संपूर्ण किंवा निर्णेय असतात या
अधिक व्यापक प्रश्नाचे---उत्तर शोधताना पुढील व्याख्याही उपयुक्त ठरेल.
संच $X$ हा रिक्त असेल किंवा $f \colon \Nat
\to X$ असे आच्छादक फलन असेल तेव्हा आणि केवळ तेव्हाच तो गणनीय
असतो, हे आठवा. अशा फलनाला $X$ चे प्रगणन म्हणतात.

\begin{defn}
संच $X$ हा रिक्त असेल किंवा त्याचे संगणनक्षम प्रगणन असेल तेव्हा आणि केवळ तेव्हाच
त्याला \emph{संगणनक्षमपणे प्रगणनीय} (संक्षेपात संगणनक्षमपणे प्रगणनीय) म्हणतात.
\end{defn}

निर्णेय स्वयंसिद्धकसंचाने निरूपणयोग्यता व्यतिरिक्त, अपूर्णतेची प्रमेये लागू होण्यासाठी
उपपत्तीविषयी आणखी एक अट अशी असेल की ती संगणनक्षम फलने आणि
निर्णेय संबंध यांच्याबद्दलची मूलभूत तथ्ये सिद्ध करू शकेल इतकी
प्रबळ असावी. ``मूलभूत तथ्ये'' म्हणजे संगणनक्षम फलनांच्या प्रत्येक
आदानासाठी त्यांची मूल्ये काय आहेत हे व्यक्त करणारी वाक्ये.
आणि ``पुरेशी प्रबळ'' म्हणजे संबंधित उपपत्ती ही वाक्ये आपली प्रमेये
मानतात. उदाहरण म्हणून बेरीज हे सर्वसाधारण संगणनक्षम फलन घ्या.
$2$ आणि $3$ या आदानांसाठी $+$ चे मूल्य~$5$ आहे, म्हणजे
$2+3 = 5$. $2$ आणि $3$ या आदानांसाठी $+$ चे मूल्य~$5$ आहे
असे सांगणारे अंकगणिताच्या भाषेतील वाक्य आहे:
$(\num{2} + \num{3}) = \num{5}$. उदाहरणार्थ, $\Th{Q}$ हे वाक्य
सिद्ध करते. अधिक सर्वसाधारणपणे, प्रत्येक संगणनक्षम फलन
$f(x_1, x_2)$ साठी~$\Lang{L_A}$ मध्ये असे सूत्र~$\varphi_f(x_1,
x_2, y)$ हवे की $f(n_1, n_2) = m$ असेल तेव्हा $\Th{Q}
\Proves \varphi_f(\num{n_1}, \num{n_2}, \num{m})$. अशा रीतीने
$\Th{Q}$ हे $n_1$, $n_2$ या आदानांसाठी~$f$ चे मूल्य~$m$
आहे हे सिद्ध करते. प्रत्यक्षात, आपण याहून थोडे अधिक मागतो:
या $n_1$,~$n_2$ या आदानांसाठी~$f$ चे मूल्य दुसरी कोणतीही संख्या नाही हेदेखील
तिने सिद्ध करावे. निर्णेय संबंधांबाबतही तसेच आहे. पुढील
दोन व्याख्या हे नेमके करतात.

\begin{defn}
सूत्र~$\varphi(x_1, \dots, x_k, y)$ हे
$f\colon \Nat^k \to \Nat$ या फलनाचे~$\Gamma$ मध्ये
\emph{निरूपण} करते तेव्हा आणि केवळ तेव्हाच की $f(n_1,
\dots, n_k) = m$ असेल तेव्हा पुढील दोन्ही अटी पूर्ण होतात:
\begin{enumerate}
\item $\Gamma \Proves \varphi(\num{n_1}, \dots, \num{n_k}, \num{m})$, आणि
\item $\Gamma \Proves \lforall[y](\varphi(\num{n_1}, \dots, \num{n_k},
y) \lif y = \num{m})$.
\end{enumerate}
\end{defn}

\begin{defn}
सूत्र~$\varphi(x_1, \dots, x_k)$ हे
$R \subseteq \Nat^k$ या संबंधाचे \emph{निरूपण} करते तेव्हा आणि केवळ तेव्हाच की
\begin{enumerate}
\item $R(n_1, \dots, n_k)$ असेल तेव्हा $\Gamma \Proves \varphi(\num{n_1},
\dots, \num{n_k})$, आणि
\item $R(n_1, \dots, n_k)$ नसेल तेव्हा $\Gamma \Proves \lnot
\varphi(\num{n_1}, \dots, \num{n_k})$.
\end{enumerate}
\end{defn}

पहिले अपूर्णता प्रमेय लागू होण्याची सामर्थ्याची अट अशी आहे: उपपत्ती
सर्व संगणनक्षम फलनांचे आणि सर्व निर्णेय संबंधांचे निरूपण
करत असावी. $\Th{Q}$ आणि तिच्या विस्तारित उपपत्ती ही अट पूर्ण करतात.
पण हे सिद्ध करण्यास आपल्याला थोडा वेळ लागेल: $\Th{Q}$ काय
सिद्ध करू शकते याविषयीचे हे सोपे नसलेले तथ्य आहे, कारण $\Th{Q}$ कडे
थोडीच स्वयंसिद्धके आहेत आणि त्यांच्यापासून ही सर्व तथ्ये सिद्ध
करावी लागतील. तरी $\Th{Q}$ ही अतिशय दुर्बल उपपत्ती आहे.
म्हणून $\Th{Q}$ सर्व संगणनक्षम फलने निरूपित करते हे सिद्ध
करणे कठीण असले, तरी बहुतेक महत्त्वाच्या उपपत्ती~$\Th{Q}$ पेक्षा
प्रबळ असतात, म्हणजे $\Th{Q}$ पेक्षा अधिक गोष्टी सिद्ध करतात.
$\Th{Q}$ ने सिद्ध केलेली गोष्ट कोणतीही अधिक प्रबळ उपपत्तीही
सिद्ध करते; म्हणून $\Th{Q}$ सर्व संगणनक्षम फलने निरूपित करत
असल्यास प्रत्येक अधिक प्रबळ उपपत्तीदेखील तसे करते. त्यामुळे
पहिल्या अपूर्णता प्रमेयाची ही अट अनेक महत्त्वाच्या उपपत्ती पूर्ण करतात.
आपल्या कठीण परिश्रमांचे फळ म्हणजे ही प्रमेये उपपत्तींच्या विस्तृत
वर्गाला लागू होतात हे दिसेल. नक्कीच, ``सर्व गणित'' आकारिकरीत्या
मांडू पाहणाऱ्या कोणत्याही उपपत्तीने $\Th{Q}$ सिद्ध करते ते
सर्व सिद्ध केले पाहिजे; किमान प्राथमिक संगणनांचे निष्कर्ष तरी
तिला सामावून घेता आले पाहिजेत. म्हणून ``सर्व गणिताची'' उपपत्ती
म्हणवण्यास पात्र असलेली उपपत्ती सुसंगत व प्रभावी स्वयंसिद्धकीय
मांडणी असलेली असेल तर तिला पहिले अपूर्णता प्रमेय लागू होईल.
दुसऱ्या अपूर्णता प्रमेयासाठी पुढील आढाव्यात सांगितलेल्या अधिक
कडक अटीही हव्यात.














\section{अपूर्णतेच्या निष्कर्षांचा आढावा}\label{inc:int:ovr:sec}

गणिताला निर्णेय स्वयंसिद्धकसंचाने निरूपणयोग्य उपपत्तीत
आकारिक रूप देता येईल आणि ती संपूर्ण व निर्णेय आहे हे
सिद्ध करता येईल, अशी हिल्बर्टची अपेक्षा होती. शिवाय, अतिशय दुर्बल,
``सांतवादी'' साधनांनी त्या उपपत्तीची सुसंगतता सिद्ध करून
अंतःप्रज्ञावादाच्या आक्षेपांपासून अभिजात गणिताचे समर्थन करण्याचे
त्याचे ध्येय होते. गोडेलच्या अपूर्णता प्रमेयांनी ही उद्दिष्टे
साध्य होऊ शकत नाहीत हे दाखवले.

गोडेलच्या पहिल्या अपूर्णता प्रमेयाने रसेल आणि व्हाइटहेड यांच्या
\emph{Principia Mathematica} ची एक आवृत्ती संपूर्ण नाही हे
दाखवले. पण सिद्धता प्रत्यक्षात अतिशय व्यापक होती आणि अनेक प्रकारच्या
उपपत्तींना लागू होते. म्हणजे केवळ \emph{Principia Mathematica}
गणित पूर्णपणे पकडू शकला नाही असे नव्हे, तर कोणतीही \emph{स्वीकारार्ह}
उपपत्ती तसे करू शकत नाही. अपूर्णता प्रमेये लागू होण्यासाठी उपपत्तीत
कोणते गुणधर्म पुरेसे आहेत हे वेगळे ओळखण्यास आणि गोडेलची सिद्धता
केवळ त्यांवर अवलंबून राहील अशी व्यापक करण्यास काही काळ लागला.
आता मात्र अंकगणिताच्या $\Lang L_A$ भाषेतील उपपत्त्यांसाठी पहिले
अपूर्णता प्रमेय अतिशय व्यापक रूपात मांडता येते.

\begin{thm}
$\Gamma$ ही~$\Lang
L_A$ मधील सुसंगत आणि निर्णेय स्वयंसिद्धकसंचाने निरूपणयोग्य उपपत्ती असेल, आणि ती सर्व
संगणनक्षम फलनांचे व निर्णेय संबंधांचे निरूपण करत असेल,
तर $\Gamma$ संपूर्ण नाही.
\end{thm}

$\Gamma$ संपूर्ण नाही याचा अर्थ किमान एक असे
वाक्य~$\varphi$ आहे की $\Gamma \Proves/ \varphi$ आणि
$\Gamma \Proves/ \lnot
\varphi$. अशा वाक्याला ($\Gamma)$ पासून \emph{स्वतंत्र}
म्हणतात. अशी स्वतंत्र वाक्ये असलीच पाहिजेत हे प्रत्यक्षात तुलनेने
लवकर सिद्ध करता येते. परंतु गोडेलच्या सिद्धतेचे सामर्थ्य असे की
ती स्वतःच्या असिद्धतायोग्यतेचा दावा करणाऱ्या वाक्याचे
\emph{विशिष्ट उदाहरण} देते.
या रंजक रचनेत~$\Gamma$ साठीचे \emph{गोडेल-वाक्य} म्हणून ओळखले
जाणारे वाक्य~$\mathsf{G}_\Gamma$ तयार होते. ते सिद्ध करता येत
नाही, कारण $\Gamma$ मध्ये $\mathsf{G}_\Gamma$ हेच वाक्य~$\Gamma$ मध्ये
$\mathsf{G}_\Gamma$ सिद्ध करता येत नाही या दाव्याशी सममूल्य असते.
ही रचना \emph{रचनात्मकपणे} केली जाते: $\Gamma$ ची स्वयंसिद्धकीय
मांडणी आणि निष्पत्ती-पद्धतीचे वर्णन दिल्यास सिद्धतेतून
$\mathsf{G}_\Gamma$ प्रत्यक्ष लिहिण्याची पद्धत मिळते. केवळ सुसंगततेवरून हे मूळ गोडेल-वाक्य
असिद्धतायोग्य असते; त्याचे खंडनही अशक्य ठरवण्यासाठी अधिक प्रबळ
गृहीतक, उदाहरणार्थ ओमेगा-सुसंगतता, आवश्यक असते. मात्र रॉसरच्या
बदललेल्या रचनेत केवळ सुसंगततेवरून स्वतंत्र वाक्य मिळते; वरील
प्रमेय त्या अधिक व्यापक रूपात दिले आहे.


गोडेलच्या सिद्धतेतील रचनेसाठी $\Lang L_A$ मधील पदे व
सूत्रे यांचे गुणधर्म आणि त्यांवरील क्रिया त्याच $\Lang L_A$
मध्ये व्यक्त करण्याचा मार्ग हवा. यात ``$\varphi$ हे वाक्य आहे'',
``$\delta$ ही~$\varphi$ ची निष्पत्ती आहे'' यांसारखे गुणधर्म
आणि $\Subst{\varphi}{t}{x}$ सारख्या क्रिया येतात. या मार्गाने
(a) चिन्हे आणि त्यांचे अनुक्रम (म्हणजे पदे, सूत्रे,
निष्पत्त्या इत्यादी) नैसर्गिक संख्यांत सांकेतिक करून
त्या गुणधर्मांचे व संबंधांचे निरूपण करावे लागते, कारण $\Lang L_A$
नैसर्गिक संख्यांविषयी बोलू शकते. तसेच (b) $\Gamma$ ला संबंधित
तथ्ये सिद्ध करता येतील अशा रीतीने हे करावे लागते. म्हणून हे
गुणधर्म नैसर्गिक संख्यांच्या निर्णेय गुणधर्मांनी आणि क्रिया
त्यांवरील संगणनक्षम फलनांनी सांकेतिक होतात हे दाखवावे लागते.
याला ``विन्यासमीमांसेचे अंकगणितीकरण'' म्हणतात.

विन्यासमीमांसेचे अंकगणितीकरण कसे करावे हे पाहण्यापूर्वी, $\Gamma$
``पुरेशी प्रबळ'' आहे या अटीचा विचार करू; म्हणजे ती सर्व संगणनक्षम
फलनांचे आणि निर्णेय संबंधांचे निरूपण करते. यासाठी
``संगणनक्षम'' याची नेमकी व्याख्या द्यावी लागेल. ती अनेक मार्गांनी
देता येते; उदाहरणार्थ, ट्यूरिंग यंत्रांच्या प्रतिमानातून किंवा
सर्वसाधारण वापराच्या संगणक-भाषेतील कार्यक्रमांनी संगणित होणाऱ्या
फलनांच्या रूपात. मात्र आपले उद्दिष्ट या फलनांचे आणि संबंधांचे
भाषा~$\Lang L_A$ मधील उपपत्तीत निरूपण करणे असल्याने
केवळ संख्यात्मक फलनांच्या संगणनक्षमतेची सोपी व्याख्या निवडणे
उत्तम. ही \emph{पुनरावर्ती फलन} ही संकल्पना आहे. म्हणून प्रथम
पुनरावर्ती फलनांची चर्चा करू. मग $\Th{Q}$ आधीच सर्व पुनरावर्ती
फलनांचे आणि संबंधांचे निरूपण करते हे दाखवू. त्यामुळे $\Th{Q}$
आणि $\Th{PA}$ यांसारख्या विशिष्ट उपपत्त्यांना अपूर्णता प्रमेय
लागू करता येईल, कारण त्या ``पुरेशा प्रबळ'' उपपत्ती असल्याचे
आपण प्रस्थापित केलेले असेल.

विन्यासमीमांसेच्या अंकगणितीकरणातून अखेरीस सूत्र
$\OProv[\Gamma](x)$ मिळते. सूत्रे ना संख्यांनी सांकेतिक
करण्यामुळे हे सूत्र~$\Gamma$ च्या स्वयंसिद्धकांपासून सिद्धतायोग्यता
व्यक्त करते. विशेषतः, $\varphi$ चा संकेतांक~$n$ असेल आणि $\Gamma \Proves
\varphi$ असेल, तर $\Gamma \Proves \Prov[\Gamma](\num{n})$.
$\Gamma$ साठीचे हे ``सिद्धतायोग्यता विधेय'' विशिष्ट अर्थाने
$\Gamma$ ची सुसंगतताही~$\Lang L_A$ मधील वाक्य म्हणून
व्यक्त करू देते. $\Gamma$ साठीचे ``सुसंगतता-विधान'' हे
$\lnot \Prov[\Gamma](\num{n})$ हे वाक्य मानू, जिथे $n$ हा
एखाद्या विरोधाभासाचा, उदा. ~ $\lfalse$ चा, संकेतांक आहे.
दुसरे अपूर्णता प्रमेय सांगते की सुसंगत निर्णेय स्वयंसिद्धकसंचाने निरूपणयोग्य उपपत्ती
स्वतःची सुसंगतता-विधानेही सिद्ध करू शकत नाहीत. हे प्रमेय लागू
होण्यासाठी लागणाऱ्या अटी, उपपत्ती सर्व संगणनक्षम फलने आणि
निर्णेय संबंध निरूपित करते एवढ्यापेक्षा काहीशा अधिक
कडक आहेत; पण $\Th{PA}$ त्या पूर्ण करते हे आपण दाखवू.












\section{अनिर्णेयता आणि अपूर्णता}\label{inc:int:dec:sec}

गोडेलच्या अपूर्णता प्रमेयांच्या सिद्धतेसाठी विन्यासमीमांसेचे
अंकगणितीकरण आवश्यक आहे. पण ते न करताही, एखादी उपपत्ती सर्व
निर्णेय संबंधांचे निरूपण करते या गृहीतकावरून काही
चांगले निष्कर्ष मिळतात. ही सिद्धता थांबण्याच्या समस्येची अनिर्णेयता
सिद्ध करणाऱ्या विकर्ण युक्तिवादासारखी आहे.

\begin{thm}
$\Gamma$ ही प्रत्येक निर्णेय संबंधांचे निरूपण
करणारी सुसंगत उपपत्ती असेल, तर $\Gamma$ निर्णेय नाही.
\end{thm}

\begin{proof}
$\Gamma$ निर्णेय आहे असे समजू. $\Gamma$ प्रत्येक
निर्णेय संबंधांचे निरूपण करत असेल तर विसंगत असलीच
पाहिजे, हे दाखवू.

निर्णेय गुणधर्म (एकस्थानी संबंध) एक मुक्त चल असलेल्या
सूत्रांनी निरूपित होतात. येथे एकच मुक्त चल असलेल्या सर्व
सुविन्यस्त सूत्रांचे प्रगणन घ्यायचे आहे; फक्त निर्णेय गुणधर्मांचे
निरूपण करणाऱ्या सूत्रांचा उपवर्ग निवडायचा नाही. त्यांचे
$\varphi_0(x)$, $\varphi_1(x)$, \dots, हे संगणनक्षम प्रगणन असू द्या.
आता पुढील संच $D \subseteq \Nat$ पाहा:
\[D = \Setabs{n}{\Gamma \Proves \lnot \varphi_n(\num{n})}
\]
$D$ हा निर्णेय संच आहे. कारण प्रथम $\varphi_n(x)$ संगणित
करून, त्यावरून $\lnot \varphi_n(\num{n})$ मिळवून $n \in D$
आहे का हे तपासता येते. $\varphi_n(x)$ मधील $x$ च्या प्रत्येक मुक्त
आस्थितीऐवजी पद $\num{n}$ ठेवणे आणि $\varphi_n(\num{n})$ च्या
आधी $\lnot$ लावणे ही यांत्रिक प्रक्रिया आहे. गृहीतकानुसार
$\Gamma$ निर्णेय असल्यामुळे
$\lnot \varphi_n(\num{n}) \in \Gamma$ आहे का हे तपासता येते.
असल्यास $n \in D$, नसल्यास $n \notin D$. त्यामुळे $D$
देखील निर्णेय आहे.

$\Gamma$ सर्व निर्णेय गुणधर्मांचे निरूपण करत
असल्याने ती~$D$ चे देखील निरूपण करते. $\Gamma$ मध्ये
$D$ चे निरूपण करणारी सर्व सूत्रे ही $\varphi_0(x)$,
$\varphi_1(x)$, \dots, या यादीत आहेत. म्हणून $\varphi_d(x)$ हे
$\Gamma$ मध्ये $D$ चे निरूपण करेल असा क्रमांक $d$
घ्या. $d
\notin D$ असेल, तर $\varphi_d(x)$ हे~$D$ चे निरूपण
करत असल्यामुळे $\Gamma \Proves
\lnot \varphi_d(\num{d})$. पण मग $d$ हा~$D$ च्या व्याख्येतील
अट पूर्ण करतो, म्हणून $d \in D$. हे $d \notin D$ शी
विसंगत आहे. म्हणून अप्रत्यक्ष सिद्धतेने $d \in D$.

$d \in D$ असल्यामुळे~$D$ च्या व्याख्येनुसार $\Gamma \Proves \lnot
\varphi_d(\num{d})$. दुसरीकडे $\varphi_d(x)$ हे $\Gamma$ मध्ये~$D$ चे
निरूपण करत असल्यामुळे $\Gamma \Proves \varphi_d(\num{d})$.
म्हणून $\Gamma$ विसंगत आहे.
\end{proof}

\begin{explain}
वरील प्रमेय दाखवते की सर्व निर्णेय संबंधांचे निरूपण
करणारी कोणतीही सुसंगत उपपत्ती निर्णेय असू शकत नाही.
$\Th{Q}$ सर्व निर्णेय संबंधांचे निरूपण करते हे
आपण दाखवू. म्हणजे $\Th{Q}$ समाविष्ट असलेल्या $\Th{PA}$
आणि $\Th{TA}$ यांसारख्या सर्व उपपत्तीदेखील तसे करतात; त्यांपैकी
सुसंगत उपपत्ती निर्णेय नसतात. (येथे नाव घेतलेल्या उपपत्ती
प्रमाण प्रतिमानात सत्य असल्यामुळे सुसंगत आहेत.)

या निष्कर्षाचा वापर करून पहिल्या अपूर्णता प्रमेयाची दुर्बल आवृत्तीही
मिळते. निर्णेय स्वयंसिद्धकसंचाने निरूपणयोग्य आणि संपूर्ण असलेली कोणतीही
उपपत्ती निर्णेय असते. त्यामुळे सुसंगत, निर्णेय स्वयंसिद्धकसंचाने निरूपणयोग्य
आणि सर्व निर्णेय गुणधर्मांचे निरूपण करणारी उपपत्ती
संपूर्ण असू शकत नाही.
\end{explain}

\begin{thm}
$\Gamma$ निर्णेय स्वयंसिद्धकसंचाने निरूपणयोग्य आणि संपूर्ण असेल, तर ती
निर्णेय असते.
\end{thm}

\begin{proof}
कोणतीही विसंगत उपपत्ती निर्णेय असते, कारण विसंगत
उपपत्त्यांत सर्व वाक्ये असतात. म्हणून ``$\varphi \in
\Gamma$ आहे का?'' या प्रश्नाचे उत्तर नेहमी ``होय'' असते,
म्हणजे तो निर्णय करता येतो.

आता $\Gamma$ सुसंगत, निर्णेय स्वयंसिद्धकसंचाने निरूपणयोग्य आणि संपूर्ण आहे
असे समजू. $\Gamma$ निर्णेय स्वयंसिद्धकसंचाने निरूपणयोग्य असल्याने ती
संगणनक्षमपणे प्रगणनीय आहे. कारण~$\Gamma$ च्या
स्वयंसिद्धकांपासूनच्या सर्व योग्य निष्पत्त्यांचे संगणनक्षम
फलनाने प्रगणन करता येते. योग्य निष्पत्ती पासून ती ज्या
वाक्यानिष्पन्न करते ते संगणित करता येते. त्यामुळे
दोन्ही मिळून~$\Gamma$ च्या सर्व प्रमेयांचे प्रगणन करणारे
संगणनक्षम फलन मिळते. $\Gamma$ सुसंगत आणि संपूर्ण असल्यामुळे
वाक्य~$\varphi$ हे~$\Gamma$ चे प्रमेय असते तेव्हा आणि केवळ तेव्हाच
$\lnot \varphi$ प्रमेय नसते. म्हणून $\varphi \in
\Gamma$ आहे का हे पुढीलप्रमाणे ठरवता येते. $\Gamma$ च्या सर्व
प्रमेयांचे प्रगणन करा. या यादीत $\varphi$ आल्यास $\Gamma \Proves \varphi$
हे कळते. यादीत $\lnot
\varphi$ आल्यास $\Gamma \Proves/ \varphi$ हे कळते. $\Gamma$
संपूर्ण असल्यामुळे यांपैकी एक प्रकरण अखेरीस घडतेच;
म्हणून प्रक्रिया अखेरीस उत्तर देते.
\end{proof}

\begin{cor}
\label{inc:int:dec:cor:incompleteness}
$\Gamma$ सुसंगत, निर्णेय स्वयंसिद्धकसंचाने निरूपणयोग्य आणि प्रत्येक निर्णेय
गुणधर्मांचे निरूपण करणारी असेल, तर ती संपूर्ण नाही.
\end{cor}

\begin{proof}
$\Gamma$ संपूर्ण असेल, तर मागील प्रमेयामुळे ती
निर्णेय असेल (कारण ती निर्णेय स्वयंसिद्धकसंचाने निरूपणयोग्य आणि सुसंगत आहे).
पण $\Gamma$ प्रत्येक निर्णेय गुणधर्मांचे निरूपण
करत असल्यामुळे पहिल्या प्रमेयाने ती निर्णेय नाही.
\end{proof}

\begin{prob}
$\Th{TA} = \Setabs{\varphi}{\Sat{N}{\varphi}}$ ही निर्णेय स्वयंसिद्धकसंचाने निरूपणयोग्य
नाही, हे दाखवा. $\Th{TA}$ सर्व निर्णेय गुणधर्मांचे निरूपण करते
असे गृहीत धरू शकता.
\end{prob}

उदाहरणार्थ, $\Th{Q}$ सर्व निर्णेय गुणधर्मांचे निरूपण
करते हे आपण प्रस्थापित केल्यावर उपप्रमेयावरून $\Th{Q}$ अपूर्ण
असलीच पाहिजे असे कळेल. पण त्याची सिद्धता स्वतंत्र वाक्य
चे उदाहरण देत नाही; असे वाक्य अस्तित्वात असलेच पाहिजे
एवढेच दाखवते. त्यासाठी विन्यासमीमांसेचे अंकगणितीकरण करून
गोडेलच्या मूळ सिद्धतेची कल्पना वापरावी लागेल. आणि अर्थातच,
$\Th{Q}$ सर्व निर्णेय गुणधर्मांचे निरूपण करते हा
पहिला दावाही अजून सिद्ध करायचा आहे.

प्रत्येक \emph{रंजक} उपपत्ती अपूर्ण किंवा अनिर्णेय असेलच असे
नाही. अनेक उपपत्ती महत्त्वाची गणितीय तथ्ये व्यक्त करण्याइतक्या
प्रबळ असूनही गोडेलच्या निष्कर्षांच्या अटी पूर्ण करत नाहीत.
उदाहरणार्थ, $\Th{Pres} = \Setabs{\varphi \in \Lang{L_{A^+}}}{\Sat{N}{\varphi}}$
ही प्रमाण प्रतिमानात सत्य असलेल्या, पण~$\times$ नसलेल्या
अंकगणिताच्या भाषेतील वाक्यांची उपपत्ती संपूर्ण आणि
निर्णेय दोन्ही आहे. या उपपत्तीला प्रेस्बर्गर अंकगणित म्हणतात;
केवळ $\Obj{0}$, $\prime$ आणि~$+$ यांच्या साहाय्याने मांडता
येणारी नैसर्गिक संख्यांविषयीची सर्व सत्ये ती सिद्ध करते.



